% !TeX root = ../main.tex
\begin{abstract}
Visual world models offer a foundation for low-altitude unmanned aerial vehicle (UAV) simulation and navigation by predicting future observations along planned flight trajectories.
However, missing camera annotations and inconsistent coordinate conventions and translation scales limit the reliability of supervision assembled from heterogeneous video sources.
For camera-conditioned models trained only on RGB targets, scene structure remains implicit in the prediction objective, without direct depth supervision to constrain viewpoint-dependent geometry.
We propose \textbf{LAVi}, a unified framework for \textbf{L}ow-\textbf{A}ltitude \textbf{Vi}sual world modeling that couples camera-annotated data construction with geometry-supervised video generation.
We assemble 600,000 clips spanning simulation, public videos, and real UAV flights, including 200,000 simulated clips whose ground-truth trajectories guide metric-scale calibration and pose-estimator selection for camera annotation across sources.
A shared transformer with modality-specific prediction branches jointly generates RGB frames and depth maps, using metric-depth supervision to constrain scene geometry and scale.
We directly adapt the bidirectionally trained model to few-step causal generation by combining future-frame alignment under reference histories with feature-space distribution matching of self-generated rollouts.
Experiments show that our data construction strategy and geometry-supervised generation method together enable LAVi to achieve the best camera-control accuracy and visual fidelity among the evaluated world models.
\end{abstract}
